\paragraph{Interpretare.} Diferențele apar la proporții mari de scrieri. La 50\% scrieri, invalidarea păstrează
doar \R{e3/invalidate/0.5/hit_ratio}\% hit-uri (fiecare scriere șterge o cheie fierbinte), iar debitul este
\R{e3/invalidate/0.5/throughput_rps} cereri/s. Write-through ridică rata de hit la
\R{e3/write-through/0.5/hit_ratio}\%. Write-behind atinge \R{e3/write-behind/0.5/throughput_rps} cereri/s
(\R{e3/write-behind/0.5/hit_ratio}\% hit-uri) și reduce numărul de tranzacții la \R{e3/write-behind/0.5/tx1k} per
1.000 de scrieri, față de $\approx$1.000 la strategiile sincrone. Câștigul vine din grupare (mai multe
modificări într-un \texttt{UPDATE}) și din combinarea modificărilor repetate ale aceluiași produs fierbinte.

La 95\% citiri strategiile converg, iar write-behind este chiar puțin mai lent. Loturile sunt mici (rata de
scriere e redusă), deci gruparea aduce puțin, iar fiecare scriere plătește un \emph{read-modify-write} în cache.

\paragraph{Lecție de proiectare: backlog-ul.} O primă versiune folosea un singur worker per instanță. Debitul
măsurat era mai mare, dar la finalul rulării erau $\approx$181.000 de scrieri încă în coadă, golite în
$\approx$5~s după oprirea încărcării. Sistemul confirma deci scrieri mai repede decât le putea persista. Am
împărțit write-behind în 4 benzi paralele, partiționate după \texttt{id mod 4}, ceea ce păstrează ordinea
modificărilor fiecărui produs. Cu ele coada rămâne goală la final, iar cifrele de mai sus reprezintă un debit
\emph{sustenabil}. Fereastra de inconsistență dintre cache și baza de date este de ordinul unui lot (milisecunde).
